\section{Introduction}\label{sec:1}

Complex Hessian equations are important examples of fully non-linear PDE's of second order on complex manifolds. They interpolate between (linear) complex \hbox{Poisson} equations ($m = 1$) and (non linear) complex Monge-Ampère equations \hbox{($m=n$)}.
They appear in many geometric problems, including the $J$-flow
\cite{SW} and quaternionic geometry \cite{AV}. They have attracted the attention of many researchers these last years as we will mention below.

\subsection{Statement of the problem}

Let $\Omega \Subset \C^n$ be a bounded domain and let $m$ be a fixed integer such that $1 \leq m \leq n$.
We consider the following general Dirichlet problem for the complex $m$-Hessian equation.

\subsubsection*{The Dirichlet problem} Let $g \in \Cc^{0} (\partial \Omega)$ be a continuous function (the boundary data) and $\mu$ be a positive Borel measure on $\Omega$ (the right hand side). The problem is to find a necessary and sufficient condition on $\mu$ such that the following problem admits a solution:
\begin{equation}\label{eq:DirPb}
\begin{cases}
U \in \mathcal{SH}_m (\Omega) \cap \mathcal {C}^{0} ({\Omega}) \\
(dd^c U)^m \wedge \beta^{n - m} = \mu &\text{on } \Omega\;\;\quad(\dag)\\
U_{\mid \partial \Omega} = g & \text{on } \partial \Omega\quad(\dag \dag)
\end{cases}
\end{equation}
The equation $(\dag)$ must be understood in the sense of currents on $\Omega$ as it will be explained in Section $2$.
The equality $(\dag \dag)$ means that $\lim_{z \to \zeta} U (z) = g (\zeta)$ for any point $\zeta \in \partial \Omega$.

Recall that for a real function $u \in \Cc^2 (\Omega)$ and each integer $1 \leq k \leq n$, we denote by $\sigma_k (u)$ the continuous function defined at each point $z \in \Omega$ as the $k$-th symmetric polynomial of the eigenvalues $\lambda (z):= (\lambda_1 (z),\dots \lambda_n(z))$
of the complex Hessian matrix $\Bigl(\frac{\partial^2 u }{\partial z_j \partial \bar{z}_k} (z)\Bigr)$ of $u$ \ie
\[
\sigma_k (u) (z):= \sum_{1 \leq j_1 < \cdots < j_k \leq n} \lambda_{j_1} (z) \cdots \lambda_{j_k} (z), \, \, \, \, z \in \Omega.
\]

We say that a real function $u \in \Cc^2 (\Omega)$ is $m$-subharmonic on $\Omega$ if for any $k$ such that $1 \leq k \leq m$, we have $\sigma_k (u) \geq 0$ pointwise on $\Omega$.

For $m= 1$, $\sigma_1 (u) = (1 \slash 4) \Delta u$ and for $m = n$, $\sigma_n (u) = \det \Bigl(\frac{\partial^2 u }{\partial z_j \partial \bar{z}_k} (z)\Bigr)$.
Therefore, $1$-subharmonic means subharmonic and $n$-subharmonic means plurisubharmonic.

As observed by Z.\,B\l ocki (\cite{Bl05}), it is possible to define a general notion of $m$-subharmonic functions using the theory of $m$-positive currents (see Section~\ref{sec:2}). Moreover, it is possible to define the $k$-Hessian measure $(dd^c u)^k \wedge \beta^{n - k}$ when \hbox{$1 \leq k \leq m$} for any (locally) bounded $m$-subharmonic function $u$ on $\Omega$ (see Section~\ref{sec:2}).

When $\mu = 0$, the Dirichlet problem \eqref{eq:DirPb} can be solved using the Perron method as for the complex Monge-Ampère equation (see \cite{Bl05}, \cite{Ch16a}).

When $g = 0$ and $\mu$ is a positive Borel measure on $\Omega$, the Dirichlet problem is much more difficult. A necessary condition for the existence of a solution to \eqref{eq:DirPb} is the existence of a subsolution.

Therefore, a particular case of the Dirichlet problem \eqref{eq:DirPb} we are interested in can be formulated as follows.

\subsubsection*{The Hölder continuous subsolution problem}
Let $\mu$ be a positive Borel measure on~$\Omega$. Assume that there exists a function $\varphi \in \mathcal{SH}_m (\Omega) \cap \Cc^{\alpha} (\bar{\Omega})$ satisfying the following condition:
\begin{equation} \label{eq:subsolution}
\mu \leq (dd^c \varphi)^m \wedge \beta^{n - m}\text{ on }\Omega, \quad \text{and}\quad \varphi_{\mid \partial \Omega} = 0.
\end{equation}

\begin{enumerate}
\item
Does the Dirichlet problem \eqref{eq:DirPb} admit a Hölder continuous solution $U_{\mu,g}$ for any boundary data $g$ which is Hölder continuous on $\partial{\Omega}$?

\item
In this case, is it possible to estimate precisely the Hölder exponent of the solution $U_{\mu,g}$ in terms of the Hölder exponents of $\varphi$ and $g$?
\end{enumerate}

Our goal in this paper is to answer the first question on the existence of a Hölder continuous solution and give an explicit lower bound of the Hölder exponent of the solution in terms of the Hölder exponent of the subsolution when the measure $\mu$ has finite total mass.

\subsection{Known results}
There have been many articles on the subject. We will only mention those that are relevant to our study and closely related to our work.
The terminology used below will be defined in the next section.

Assume that $\Omega$ is a smooth strongly $m$-pseudoconvex domain. When the boundary data $g$ is smooth and the right hand side $\mu = f \lambda_{2n}$ is a measure with a smooth positive density $f > 0$, S.\,Y.\,Li proved in \cite{Li04} that the problem has a unique smooth solution. Later, Z.\,B\l ocki introduced the notion of weak solution and solved the Dirichlet problem for the homogeneous Hessian equation in the unit ball in $\C^n$ (\cite{Bl05}).
When the density $0 \leq f \in L^p (\Omega)$ with $p > n \slash m$, Dinew and Ko\l odziej proved the existence of a continuous solution (\cite{DK14}). Assuming moreover that $g$ is Hölder continuous on~$\bar \Omega$, Ngoc Cuong Nguyen proved the Hölder continuity of the solution under an additional assumption on the density $f$ (\cite{N14}) and M.\,Charabati proved the Hölder continuity of the solution for general densities (\cite{Ch16}).

On the other hand, S.\,Ko\l odziej \cite{Kol05} proved that the Dirichlet problem has a bounded plurisubharmonic solution if (and only if) it has a bounded subsolution with zero boundary values. This is known as the bounded subsolution theorem for plurisubharmonic functions. The same result was proved for the Hessian equation by Ngoc Cuong Nguyen in \cite{N13}.

The Hölder continuous subsolution problem stated above has attracted a lot of attention these last years and was formulated in \cite{DGZ16} for the complex Monge-Ampère equation.

\smallskip
It has been solved for the complex Monge-Ampère by Ngoc Cuong Nguyen in \cite{N18a,N18b}.
Recently S.\,Ko{\l}odziej and Ngoc Cuong Nguyen solved the Hölder subsolution problem for the Hessian equation under the restrictive assumption that the measure $\mu$ is compactly supported on $\Omega$ (see \cite{KN18}, \cite{KN19}).

\subsection{Main new results}

In this paper we will solve the Hölder continuous subsolution problem for Hessian equations when $\mu$ is any positive Borel measure with finite mass on $\Omega$.

Our first main result gives a new comparison inequality which will be applied to positive Borel measures without restriction on their support.

\begin{theo}\label{th:A}
Let $\Omega \Subset \C^n$ be a bounded strongly $m$-pseudoconvex domain.
Let $\varphi\in \mathcal{SH}_m(\Omega)\cap \Cc^{\alpha}(\overline\Omega)$ with $0 < \alpha \leq 1$ such that $\varphi = 0$ in $\partial \Omega$. Then for any $r$ such that $0 < r < m \slash (n-m)$, there exists a constant $A>0$ such that for every compact $K\subset\Omega$,
\[
\int_ K (dd^c\varphi)^m\wedge\beta^{n-m} \leq A\Bigl(\left[\Capa_m (K,\Omega)\right]^{1 + \epsilon} + \left[\Capa_m (K,\Omega)\right]^{1 + m\epsilon}\Bigr),
\]
where $\epsilon:= \spfrac{\alpha r}{(2-\alpha) m + \alpha} > 0$.
\end{theo}

The capacity $\Capa_m (K,\Omega)$ will be defined in the next section. The constant $A$ in the theorem is explicit (see Formula \eqref{eq:finalConst}). Observe that the most relevant case in the application of this inequality will be when $\Capa_m (K,\Omega) \leq 1$. In this case the right exponent is $\tau = 1 + \spfrac{\alpha r}{(2-\alpha) m + \alpha}$.

Theorem~\ref{th:A} substantially improves a recent result of \cite{KN19} who proved an estimate of this kind when the compact set $K \subset \Omega'$ is contained in a fixed open set $\Omega' \Subset \Omega$, \ie $K$ stays away from the boundary of $\Omega$.

When $m=n$ a better estimate was obtained in \cite{N18a} using the exponential integrability of plurisubharmonic functions which fails when $m < n$.

As a consequence of Theorem~\ref{th:A}, we will deduce the following result which solves the Hölder continuous subsolution problem.

\begin{theo}\label{th:B}
Let $\Omega \Subset \C^n$ be a bounded strongly $m$-pseudoconvex domain and $\mu $ a positive Borel measure on $\Omega$ with finite mass. Assume that there exists $\varphi\in \mathcal{E}^0_m(\Omega)\cap\Cc^{\alpha}(\overline\Omega)$ with $0 < \alpha < 1$ such that
\begin{equation} \label{eq:subsol}
\mu \leq (dd^c\varphi)^m\wedge\beta^{n-m} \text{ weakly on }\Omega, \quad \text{and}\quad \varphi \equiv 0 \text{ on }\partial \Omega.
\end{equation}
Then for any continuous function $g \in \Cc^{2 \alpha} (\partial \Omega)$, there exists a unique function $U = U_{g,\mu} \in \mathcal{SH}_m (\Omega) \cap \Cc^{0} (\bar{\Omega})$ such that
\[
(dd^c U)^m\wedge\beta^{n-m} = \mu, \quad\text{and}\quad U = g \text{ on } \partial \Omega.
\]
Moreover, $U \in \Cc^{\alpha'}(\overline\Omega)$ for any $\alpha'$ such that $0 < \alpha' < \Psfrac{\alpha^m}{2^{m - 1}}\cdot \gamma (m,n,\alpha)$, where
\[
\gamma (m,n,\alpha):= \frac{m \alpha}{ m (m + 1) \alpha + 2 (n-m)}\cdot
\]
\end{theo}

Recall that by definition when $\alpha = 1 \slash 2$, $g \in \Cc^1 (\partial \Omega)$ means that $g$ is Lipschitz and when $1 \slash 2 < \alpha <1$ and $2 \alpha = 1 + \theta$ with $0 < \theta < 1$, $ g \in \Cc^{2 \alpha} (\partial \Omega)$ means that $g \in \Cc^1 (\partial \Omega)$ and and $\nabla g$ is Hölder continuous of exponent $\theta$ on $\partial \Omega$.

Let us give a rough idea of the proofs of these results.

\subsubsection*{Idea of the proof of Theorem~\ref{th:A}}
The general idea of the proof is inspired by \cite{KN19}. However, since our measure is not compactly supported nor of finite mass, we need to control the behaviour of the $m$-Hessian measure of $\varphi$ close to the boundary. This will be done in several steps in Section~\ref{sec:3} and Section~\ref{sec:4}.
\begin{itemize}
\item
The first step is to estimate the mass of the $m$-Hessian measure $\sigma_m (\varphi)$ of a Hölder continuous $m$-subharmonic function $\varphi$ in terms of its regularization $\varphi_\delta$ on any compact set in $\Omega_\delta$. This requires to consider the $m$-subharmonic envelope of $\varphi_\delta$ on $\Omega$ and provide a precise control on its $m$-Hessian measure (see Theorem \ref{thm:obstacle}).

\item
The second step is to estimate the mass of $\sigma_m (\varphi)$ on a compact set close to the boundary in terms of its Hausdorff distance to the boundary (see Lemma \ref{lem:ComparisonIneq}).
\end{itemize}

\subsubsection*{Idea of the proof of Theorem~\ref{th:B}}
The proof will be in two steps.
\begin{itemize}
\item
The first step relies on a standard method which goes back to \cite{EGZ09} (see also \cite{GKZ08}) in the case of the complex Monge-Ampère equation. This method consists in proving a semi-stability inequality estimating
$ \sup_{\Omega} (v-u)_+ $ in terms of $\Vert (v-u)_+\Vert_{L^{1} (\Omega,\mu)}$, where $u $ is the bounded $m$-subharmonic solution to the Dirichlet problem \eqref{eq:DirPb} and $v$ is any bounded $m$-subharmonic function with the same boundary values as $u$, under the assumption that the measure $\mu$ is dominated by the $m$-Hessian capacity with an exponent $\tau > 1$ (see Definition \ref{def:cap-domination}).

\item
The second step uses an idea which goes back to \cite{DDGKPZ15} in the setting of compact Kähler manifolds (see also \cite{GZ17}). It has been also used in the local setting in \cite{N18a} and \cite{KN19}. It consists in estimating the $L^1 (\mu)$-norm of $v - u$ in terms of the $L^1 (\lambda_{2 n})$-norm of $(v-u)$ where $u$ is the bounded solution to the Dirichlet problem and $v$ is a bounded $m$-subharmonic function on $\Omega$ close to the regularization $ u_\delta$ of $u$. This step requires that the measure $\mu$ is well dominated by the $m$-Hessian capacity, which is precisely the content of our Theorem~\ref{th:A}.
Then using the Poisson-Jensen formula as in \cite{GKZ08}, we see that the $L^1$-norm of $(u_\delta-u)$ is $O (\delta^2)$ (see Lemma \ref{lem:Poisson-Jensen}) and Lemma \ref{lem:sup-mean} allows us to finish the proof.
\end{itemize}

\subsubsection*{Acknowledgements}
The authors are indebted to Hoang Chinh Lu for his very careful reading of the first version of this paper and for valuable comments that helped to correct some errors and to improve the presentation of the paper. They also would like to thank Vincent Guedj for interesting discussions and Ngoc Cuong Nguyen for useful exchanges about his earlier work on this subject.

This project started when the first author was visiting the Institute of Mathematics of Toulouse (IMT) during the springs 2017 and 2018. She would like to thank IMT for the invitation and for providing excellent research conditions.

The authors thank the referees for their useful comments.

\section{Preliminary results}\label{sec:2}
In this section, we recall the basic properties of $m$-subharmonic functions and some results we will use throughout the paper.

\subsection{Hessian potentials}
For a hermitian $n \times n$ matrix $a = (a_{j,\bar k})$ with complex coefficients, we denote by $\lambda_1,\dots \lambda_n$ the eigenvalues of the matrix $a$. For any $k$ such that $1 \leq k \leq n$ we define the $k$-th trace of $a$ by the formula
\[
s_k (a):= \sum_{1 \leq j_1 < \cdots < j_k \leq n} \lambda_{j_1} \cdots \lambda_{j_k},
\]
which is the $k$-th elementary symmetric polynomial of the eigenvalues $(\lambda_1,\dots, \lambda_n)$ of~$a$.

Let $\C^n_{(1,1)} $ be the space of real $(1, 1)$-forms on $\C^n$ with constant
coefficients, and define the cone of $m$-positive $(1,1)$-forms on $\C^n$ by
\begin{equation}\label{eq:mpositive}
\Theta_m:= \{\theta \in \C^n_{(1,1)} \sep \theta \wedge \beta^{n - 1} \geq 0,\dots, \theta^m \wedge \beta^{n - m} \geq 0\}.
\end{equation}

\skpt
\begin{definition} \label{def:mpositive}
\begin{enumerate}
\item
A smooth $(1,1)$-form $\theta$ on $\Omega$ is said to be $m$-positive on $\Omega$ if for any $z \in \Omega$, $\theta (z) \in \Theta_m$.

\item
A function $u:\Omega \to \R\cup\{-\infty\}$ is said to be $m$-subharmonic on $\Omega$ if it is subharmonic on $\Omega$ (not identically $-\infty$ on any component) and for any collection of smooth $m$-positive $(1,1)-$forms $\theta_1,\dots,\theta_{m-1}$ on $\Omega$, the following inequality
\[
dd^c u\wedge \theta_1\wedge\cdots\wedge\theta_{m-1} \wedge \beta^{n-m}\geq 0,
\]
holds in the sense of currents on $\Omega$.
\end{enumerate}
\end{definition}

We denote by $\mathcal{SH}_m (\Omega) $ the positive convex cone of $m$-subharmonic functions on~$\Omega$. We give below the most basic properties of $m$-subharmonic functions that will be used in the sequel.

\skpt
\begin{proposition}\label{prop:basic}
\begin{enumerate}
\item
If $u\in \Cc^2(\Omega)$, then $u$ is $m$-subharmonic on $\Omega$ if and only if $(dd^c u)^k\wedge \beta^{n-k}\geq0$
pointwise on $\Omega$ for $k=1,\dots, m$.

\item
$\mathcal{PSH}(\Omega)=\mathcal{SH}_n(\Omega)\subsetneq \mathcal{SH}_{n-1}(\Omega)\subsetneq\cdots\subsetneq \mathcal{SH}_1(\Omega)=\mathcal{SH}(\Omega) $.

\item
$\mathcal{SH}_m(\Omega) \subset L^1_{\loc} (\Omega)$ is a positive convex cone.

\item
If $u$ is $m$-subharmonic on $\Omega$ and $f: I \to\R$ is a convex, increasing function on some interval containing the image of $u$, then $f\circ u$ is $m$-subharmonic on $\Omega$.

\item
The limit of a decreasing sequence of functions in $\mathcal{SH}_m(\Omega)$ is $m$-subharmonic on $\Omega$ when it is not identically $- \infty$ on any component.

\item
Let $u \in \mathcal{SH}_m(\Omega)$ and $v \in \mathcal{SH}_m(\Omega') $, where $\Omega'\subset\C^n$ is an open set such that $\Omega \cap \Omega' \neq \emptyset$. If $u\geq v$ on $\Omega \cap \partial\Omega'$, then the function
\[
z \mto w(z):=\begin{cases}
\max(u(z),v(z)) &\text{if } z \in \Omega \cap \Omega'\\
u(z) &\hbox{if } z \in\Omega\setminus\Omega'
\end{cases}
\]
is $m$-subharmonic on $\Omega$.
\end{enumerate}
\end{proposition}

Another ingredient which will be important is the regularization process. Let $\chi$ be a fixed smooth positive radial function with compact support in the unit ball $\B \subset \C^n$ and $\int_{\C^n}\chi (\zeta)d\lambda_{2 n}(\zeta)=1$.
For any $\delta$ such that $ 0 < \delta < \delta_0:= \diam (\Omega)$, we set
\[
\chi_{\delta}(\zeta)=\frac{1}{\delta^{2n}}\,\chi (\sfrac{\zeta}{\delta})\quad\text{and}\quad \Omega_{\delta}=\{z \in\Omega\sep \dist (z,\partial\Omega)>\delta\}.
\]

Let $u \in \mathcal{SH}_m (\Omega) \subset L^1_{\loc} (\Omega)$ and define its standard $\delta$-regularization by the formula%
\begin{equation} \label{eq:reg}
{ u}_{\delta} (z):= \int_{\Omega} u (z - \zeta) \chi_{\delta} (\zeta) d \lambda_{2n} (\zeta), z \in \Omega_{\delta}.
\end{equation}
Then it is easy to see that $ {u}_{\delta}$ is $m$-subharmonic and smooth on $\Omega_{\delta}$ and decreases to~$u$ on $\Omega$ as $\delta $ decreases to $0$.

The following lemma was proved in \cite{GKZ08}.

\begin{lemma} \label{lem:Poisson-Jensen}
Let $u \in \mathcal{SH}_m (\Omega) \cap L^{1} (\Omega)$. Then for $\delta\in{}]0,\delta_0[$, its $\delta$-regularization extends to $\C^n$ by the formula
\begin{equation} \label{eq:regularization}
{ u}_{\delta} (z):= \int_{\Omega} u (\zeta) \chi_{\delta} (z - \zeta) d \lambda_{2n} (\zeta), z \in \C^n,
\end{equation}
and have the following properties:
\begin{enumerate}
\item
$ { u}_{\delta}$ is a smooth function on $\C^n$ which is $m$-subharmonic on $\Omega_\delta$;

\item
$({u}_\delta)$ decreases to $u$ on $\Omega$ as $\delta $ decreases to $0$ and
\[
\int_{\Omega_\delta} ( { u}_{\delta} (z) - u(z)) d\lambda_{2 n}(z) \leq a_n \delta^2 \int_{\Omega_\delta} dd^c u \wedge \beta^{n - 1},
\]
where $a_n > 0$ is a uniform constant independent of $u$ and $\delta$.
\end{enumerate}
\end{lemma}

Let us introduce the notion of strong $m$-pseudoconvexity that will be used in the sequel.

\begin{definition} We say that the open set $\Omega$ is strongly $m$-pseudoconvex if $\Omega$ admits a defining function $\rho$ which is smooth strictly $m$-subharmonic in a neighbourhood of~$\bar \Omega$ and $\vert \nabla \rho\vert > 0$ on $\partial \Omega = \{\rho = 0\}$. In this case we can choose $\rho $ so that
\begin{equation} \label{eq:stronpconvexity}
(dd^c \rho)^k \wedge \beta^{n - k} \geq \beta^n \quad \text{for $k$ such that }1 \leq k \leq m,
\end{equation}
pointwise on $\Omega$.
\end{definition}

The following lemma is analogous to a lemma proved in \cite{GKZ08} using mean values rather than convolution.
\begin{lemma} \label{lem:sup-mean} Let $\Omega \Subset \C^n$ be a bounded domain and $u \in \mathcal{SH} (\Omega) \cap L^{\infty} ({\bar \Omega})$. Assume that $u$ is Hölder continuous near $\partial\Omega$ with exponent $\alpha \in{} ]0,1[$. Then the following properties are equivalent:
\begin{enumeratei}
\item
$\exists c_1 >0$, ${ u}_\delta:= u \star \chi_\delta \leq u + c_1 \delta^{\alpha}$ in $\Omega_\delta$,

\item
$\exists c_2 >0$, $\sup_{\bar B(z,\delta)} u \leq u + c_2 \delta^{\alpha}$ in $\Omega_\delta$.
\end{enumeratei}
\end{lemma}
A similar lemma has been recently proved in the compact Hermitian manifold setting in \cite{LPT20}.
A slight modification of the proof of \cite{GKZ08} with an observation from \cite{LPT20} works also in our context as it is explained in \cite{Ze20}.

\begin{remark} \label{rem:HolderBoundary} Recall that $u$ is Hölder continuous near $\partial\Omega$ with exponent $\alpha \in{}]0,1]$ if there exists $\delta_1 > 0$ small enough and a constant $\kappa > 0$ such that for any $\zeta \in \partial \Omega$ and any $\delta\in{}]0,\delta_1[$,
\[
\sup_{z \in \Omega (\zeta,\delta)} \vert u (z) - u (\zeta) \vert \leq \kappa \delta^\alpha, \quad\text{where}\quad \Omega(\zeta,\delta):= \Omega \cap B (\zeta,\delta).
\]

Assume that there exists two functions $v, w $ defined and Hölder continuous with exponent $\alpha $ on a neighbourhood $ U$ of $\partial \Omega$ in $\bar \Omega$ such that
$v \leq u \leq w$ on $ U$ and $v = u = w$ on $\partial \Omega$. Then $u$ is Hölder continuous with exponent $\alpha $ near $\partial \Omega$.
\end{remark}

\subsection{Complex Hessian operators}

Following \cite{Bl05}, we can define the Hessian operators acting on (locally) bounded $m$-subharmonic functions as follows. Given $u_1,\dots, u_k \in \mathcal{SH}_m (\Omega) \cap L^{\infty} (\Omega)$ ($1 \leq k \leq m$), one can define inductively the following positive $(m-k,m-k)$-current on $\Omega$
\[
dd^c u_1 \wedge \cdots \wedge dd^c u_k \wedge \beta^{n - m}:= dd^c (u_1 dd^c u_2 \wedge \cdots \wedge dd^c u_k \wedge \beta^{n - m}).
\]
In particular, if $u \in \mathcal{SH}_m (\Omega) \cap L^{\infty}_{\loc} (\Omega)$, the positive current $(dd^c u)^m \wedge \beta^{n-m}$ can be identified to a positive Borel measure on $\Omega$, the so called $m$-Hessian measure of $u$ denoted by:
\[
\sigma_m (u):= (dd^c u)^m \wedge \beta^{n-m}.
\]
Observe that when $m\!=\! 1$, $\sigma_1 (u) \!=\! dd^c u \wedge \beta^{n-1}$ is the Riesz measure of $u$ (up to a positive constant), while $\sigma_n (u) \!=\! (dd^c u)^n$ is the complex Monge-Ampère measure of~$u$.

It is then possible to extend Bedford-Taylor theory to this context.
In particular, Chern-Levine Nirenberg inequalities holds and the Hessian operators are continuous under local uniform convergence and pointwise a.e. monotone convergence
on $\Omega$ of sequences of functions in $\mathcal{SH} (\Omega) \cap L^{\infty}_{\loc} (\Omega)$ (see \cite{Bl05}, \cite{Lu12}).

We define $\mathcal{E}_m^0 (\Omega) $ to be the positive convex cone of negative functions $\phi \in \mathcal{SH}^-_m (\Omega) \cap L^{\infty} (\Omega)$ with zero boundary values such that
\[
\int_{\Omega} (dd^c \phi)^m \wedge \beta^{n - m} < + \infty.
\]
These are the ``test functions'' in $m$-Hessian potential theory. The formula of integration by parts is valid for these functions.

More generally it follows from \cite{Lu12,Lu15} that the following property holds: if~$\phi \in \mathcal{E}_m^0 (\Omega) $ and $u, v \in \mathcal{SH}_m (\Omega) \cap L^{\infty} (\Omega)$ with $u \leq 0$, then for $k$ such that $0 \leq k \leq m - 1$,
\begin{equation} \label{eq:testinequality}
\int_\Omega (-\phi) dd^c u \wedge (dd^c v)^k \wedge \beta^{n - k-1} \leq \int_\Omega (-u) dd^c \phi \wedge (dd^c v)^k \wedge \beta^{n - k-1}.
\end{equation}
An important tool in the corresponding potential theory is the comparison principle.

\begin{proposition} \label{prop:comparison principle}
Assume that $u,v\in \mathcal{SH}_m(\Omega)\cap L^{\infty}(\Omega)$ and for any $\zeta \in \partial \Omega$, $\liminf_{z \to \zeta }(u(z)- v(z))\geq 0$. Then
\[
\int_{\{u<v\}}(dd^c v)^m\wedge\beta^{n-m} \leq \int_{\{u<v\}}(dd^c u)^m\wedge\beta^{n-m}.
\]
Consequently, if $(dd^cu)^m\wedge\beta^{n-m}\leq(dd^cv)^m\wedge\beta^{n-m}$ weakly on $\Omega$, then $u \geq v$ on $\Omega$.

\end{proposition}
It follows from the comparison principle that if the Dirichlet problem \eqref{eq:DirPb} admits a solution, then it is unique.

Let us recall the following estimates due to Cegrell (\cite{Ceg04}) for the complex Monge-Ampère operators and extended by Charabati to complex Hessian operators (\cite{Ch16}).

\begin{lemma} \label{lem:Cegrell} Let $u, v, w \in\mathcal{E}_m^0(\Omega)$. Then for any $1 \leq k \leq m - 1$
\[
\int_{\Omega}dd^cu\wedge(dd^cv)^k\wedge(dd^cw)^{m-k-1}\wedge\beta^{n-m}
\leq I_m (u)^{\sfrac{1}{m}} \, I_m (v)^{\sfrac{k}{m}} \, I_m (w)^{\psfrac{m-k-1}{m}},
\]
where $I_m (u):= \int_{\Omega}(dd^c u)^m \wedge \beta^{n-m}$.

In particular, if $\Omega$ is strongly $m$-pseudoconvex, then
\[
\int_{\Omega}dd^c u \wedge (dd^c w)^k \wedge\beta^{n-k -1} \leq c_{m,n} \left(I_m (u)\right)^{\sfrac{1}{m}} \left(I_m (w)\right)^{\sfrac{k}{m}},
\]
and
\[
\int_{\Omega}dd^c u \wedge \beta^{n-1} \leq c_{m,n} \left(I_m (u)\right)^{\sfrac{1}{m}},
\]
where $c_{m,n} > 0$ is a uniform constant.
\end{lemma}

The following consequence will be useful in the sequel. This result is usually stated for plurisubharmonic functions on a bounded domain with boundary values $0$. Let us give a more general version using Cegrell inequalities.

\begin{corollary} \label{coro:comparison principle} Let $\Omega \Subset \C^n$ be a bounded strongly $m$-pseudoconvex domain. Assume that $u,v\in \mathcal{SH}_m(\Omega)\cap L^{\infty}(\Omega)$ satisfy $u \leq v$ on $\Omega$ and for any $\zeta \in \partial \Omega$, $\lim_{z \to \zeta }(u(z)- v(z))= 0$.
Then
\[
\int_{\Omega}(dd^c v)^m\wedge\beta^{n-m} \leq \int_{\Omega}(dd^c u)^m\wedge\beta^{n-m}.
\]
\end{corollary}

\begin{proof} The proof is standard but let us repeat it here for the convenience of the reader.
We can assume that $I_m (u):= \int_{\Omega}(dd^c u)^m\wedge\beta^{n-m} < + \infty$.

Let $\rho: \Omega \to{} ]- \infty,0[$ be a defining $m$-subharmonic function on a neighbourhood of~$\bar \Omega$. In particular $\int_{\Omega}(dd^c \rho)^m\wedge\beta^{n-m} < + \infty$.
For fixed $\varepsilon > 0$, the function $u_\varepsilon:= u + \varepsilon \rho$ is a bounded $m$-subharmonic function such that $u_\varepsilon = v$ on $\partial \Omega$ and $\{u_\varepsilon < v\} = \Omega$.

Applying Proposition \ref{prop:comparison principle}, we obtain
\[
\int_{\Omega}(dd^c v)^m\wedge\beta^{n-m} \leq \int_{\Omega}(dd^c u_\varepsilon)^m\wedge\beta^{n-m}.
\]
Observe that
\[
(dd^c u_\varepsilon)^m\wedge\beta^{n-m} = (dd^c u)^m\wedge\beta^{n-m} + \sum_{j = 1}^m C_m^j \varepsilon^j (dd^c u)^{m - j} \wedge (dd^c \rho)^j \wedge \beta^{n -m}.
\]
By Lemma \ref{lem:Cegrell}, there exists a constant $C > 0 $ depending on $I_m (u)$ and $I_m (\rho)$ such that for any $1 \leq j \leq m$, we have
\[
\int_\Omega (dd^c u)^{m - j} \wedge (dd^c \rho)^j \wedge \beta^{n -m} \leq C.
\]
Therefore, for any $\epsilon\in{}]0,1[$,
\[
\int_{\Omega}(dd^c v)^m\wedge\beta^{n-m} \leq \int_{\Omega}(dd^c u)^m\wedge\beta^{n-m} + C 2^m \varepsilon.
\]
Letting $\varepsilon \to 0$, we obtain the required inequality.
\end{proof}

\subsection{The bounded subsolution theorem}

Let $\Omega \Subset \C^n$ be a bounded strongly $m$\nobreakdash-pseudoconvex domain. Assume there exists $v \in \mathcal{SH}_m (\Omega)\cap L^{\infty} (\Omega)$ such that
\begin{equation} \label{eq:boundedsubsol}
\mu \leq (dd^c v)^m \wedge \beta^{n -m} \text{ on } \Omega \quad\text{and}\quad v|_{\partial \Omega} \equiv 0.
\end{equation}
Ngoc Cuong Nguyen proved that under this condition, the Dirichlet problem \eqref{eq:DirPb} admits a unique bounded $m$-subharmonic solution (see \cite{N13}).

\begin{theorem}[\cite{N13}] \label{thm:boundedsubsolution}Let $\Omega \Subset \C^n$ be a bounded strongly $m$-pseudoconvex domain and $\mu$ a positive Borel measure on $\Omega$ satisfying the condition \eqref{eq:boundedsubsol}. Then for any $g \in \Cc^0 (\partial \Omega)$, there exists a unique $U = U_{g,\mu} \in \mathcal{SH}_m (\Omega)\cap L^{\infty} (\Omega)$ such that $(dd^c U)^m \wedge \beta^{n -m} = \mu$ on $\Omega$ and $U|_{\partial \Omega} \equiv g$.
\end{theorem}

\subsection{The viscosity comparison principle}

In order to prove Theorem~\ref{th:A}, we will need to prove an important result (Theorem \ref{thm:obstacle}). The proof of this result uses the viscosity comparison principle which was established for complex Hessian equations by H.\,C.\,Lu (\cite{Lu13}) in the spirit of the earlier work by P.\,Eyssidieux, V.\,Guedj and the second author on complex Monge-Ampère equations (\cite{EGZ11}).

To state this comparison principle we need some definitions. Let $\Omega \Subset \C^n$ be a bounded domain and $F: \Omega \times \R \to \R$ a continuous function \emph{non-decreasing} in the last variable.
\begin{definition}
Let $u: \Omega\to \R\cup\{-\infty\}$ be a function and $q$ be a $\Cc^2$ function in a neighborhood of $z_0\in \Omega$. We say that $q$ touches $u$ from above (\resp below) at $z_0$ if $q(z_0)=u(z_0)$ and $q(z)\geq u(z)$ (\resp $q(z)\leq u(z)$) for every $z$ in a neighborhood of $z_0$.
\end{definition}

\begin{definition}\label{def: viscosity subsolution}
An upper semicontinuous function $u: \Omega\to \R$ is a viscosity subsolution to the equation
\begin{equation}\label{eq: heq 1}
(dd^c u)^m\wedge \beta^{n-m} = F(z,u)\beta^n,
\end{equation}
if for any $z_0\in \Omega$ and any $\Cc^2$ function $q$ which touches $u$ from above at $z_0$ then
\[
\sigma_m(q) \geq F(\cdot,q (z_0))\beta^n, \quad \text{at} \ z = z_0.
\]
We will also say that $\sigma_m(u)\geq F(\cdot,u)\beta^n$ in the viscosity sense at $z_0$ and $q$ is an upper test function for $u$ at $z_0$.
\end{definition}
\begin{definition}\label{def: viscosity supersolution}
A lower semicontinuous function $v: \Omega \to \R$ is a viscosity supersolution to \eqref{eq: heq 1} if for any $z_0\in X$ and any $\Cc^2$ function $q$ which touches $v$ from below at $z_0$,
\[
[(dd^c q)^m\wedge \beta^{n-m}]_+\leq F(z,q)\beta^n, \quad \text{at} \ z = z_0.
\]
Here $[\alpha^m\wedge\beta^{n-m}]_+$ is defined to be $\alpha^m\wedge\beta^{n-m} $ if $\alpha$ is $m$-positive and $0$ otherwise. We will also say that $\sigma_m(v)_+ \leq F(\cdot,v)\beta^n$ in the viscosity sense at $z_0$ and $q$ is a lower test function for $v$ at $z_0$.
\end{definition}

\begin{remark} If $v \in \Cc^2 (\Omega)$ then $\sigma_m(v)\geq F(z,v)\beta^n$ (\resp $[\sigma_m(v)]_+\leq F(z,v)\beta^n$) holds on $\Omega$ in the viscosity sense iff it holds in the usual sense.
\end{remark}

\begin{definition}
A continuous function $u: \Omega \to \R$ is a viscosity solution to \eqref{eq: heq 1} if it is both a subsolution and a supersolution.
\end{definition}

The first important result in this theory compares the viscosity and potential subsolutions.

\begin{proposition}[\cite{Lu13}] \label{prop: viscosity vs potential general case}
Let $u$ be a bounded upper semi-continuous function in~$\Omega$. Then the inequality
\begin{equation}\label{eq: viscosity vs potential 2}
\sigma_m (u)\geq F(\cdot,u)\beta^n
\end{equation}
holds in the viscosity sense on $\Omega$ if and only if $u$ is $m$-subharmonic and \eqref{eq: viscosity vs potential 2} holds in the potential sense on $\Omega$.
\end{proposition}

Now we can state the viscosity comparison principle.

\begin{theorem} [\cite{Lu13}]\label{thm: viscosity comparison principle}
Let $u: \Omega \to \R$ be a bounded viscosity subsolution and $v: \Omega \to \R$ be a viscosity supersolution of the equation
\[
\sigma_m(u)=F(\cdot,u)\beta^n,
\]
on $\Omega$. If $u \leq v$ on $\partial \Omega$ then $u\leq v$ on $\Omega$.
\end{theorem}

For more details on this theory we refer to \cite{Lu13} and \cite{EGZ11} in the complex case and to \cite{CIL92} for the real case.

\subsection{Weak stability estimates}

An important tool in dealing with our problems is the notion of capacity. This was introduced by Bedford and Taylor in their pioneering work for the complex Monge-Ampère operator (see \cite{BT82}).
Let us recall the corresponding notion of capacity we will use here (see \cite{Lu12}, \cite{SA13}). Let $\Omega \Subset \C^n$ be a strongly $m$-pseudoconvex domain. The $m$-Hessian capacity is defined as follows. For any compact set $K \subset \Omega$,
\[
\Capa_m(K,\Omega):= \sup \bigl\{\textstyle\int_K (dd^c u)^m \wedge \beta^{n - m} \sep u \in \mathcal{SH}_m (\Omega),\, - 1 \leq u \leq 0\bigr\}.
\]

We can extend this capacity as an outer capacity on $\Omega$. Given a set $S \subset \Omega$, we define the inner capacity of $S$ by the formula
\[
\Capa_m(S,\Omega):= \sup \{\Capa_m(K,\Omega) \sep K \, \, \hbox{compact}, \, K \subset S\}.
\]
The outer capacity of $S$ is defined by the formula
\[
\Capa^*_m(S,\Omega):= \inf \{\Capa_m(U,\Omega) \sep U\text{ is open}, \,U \supset S\},
\]
It is possible to show that $\Capa^*_m(\cdot,\Omega)$ is a Choquet capacity and then any Borel set $ B \subset \Omega$ is capacitable and
for any compact set $K \subset \Omega$,
\begin{equation} \label{eq:cap}
\Capa_m(K,\Omega)=\int_{\Omega}(dd^c u_K^*)^m\wedge\beta^{n-m},
\end{equation}
where $u_K$ is the relative equilibrium potential of $(K,\Omega)$ defined by the formula:
\[
u_K:=\sup\{u\in \mathcal{SH}_m(\Omega) \sep u \, \leq \, -\mathbf{1}_K \text{ on }\Omega\},
\]
and $u_K^*$ is its upper semi-continuous regularization on $\Omega$ (see \cite{Lu12}).

It is well known that $u_K^*$ is $m$-subharmonic on $\Omega$, $- 1 \leq u_K^* \leq 0$, $u_K^*= - 1$ almost everywhere (with respect to $\Capa_m$) on $\Omega$ and $u_K^* \to 0$ as $z \to \partial \Omega$ (see \cite{Lu12}).

We will use the following definition.
\begin{definition} \label{def:cap-domination} Let $\mu$ be a positive Borel measure on $\Omega$ and let $A,\tau > 0$ be positive numbers. We say that $\mu$ is dominated by the $m$-Hessian capacity with parameters $(A,\tau)$ if for any compact subset $K \subset \Omega$ with $\Capa_m (K,\Omega) \leq 1$,
\begin{equation} \label{eq:capdomination}
\mu (K) \leq A \Capa_m (K,\Omega)^{\tau}.
\end{equation}
\end{definition}
Observe that by capacitability, this inequality is then satisfied for any Borel set $K \subset \Omega$.

Let us mention that S.\,Ko\l odziej was the first to relate the domination of the measure $\mu$ by the Monge-Ampère capacity to the regularity of the solution to complex Monge-Ampère equations (see \cite{Kol96}).

Using his idea, Eyssidieux, Guedj and the second author were able to establish in \cite{EGZ09} a weak stability $L^1$-$L^{\infty}$ estimate for bounded solutions to the Dirichlet problem for the complex Monge-Ampère equation. This result is the main tool in deriving estimates on the modulus of continuity of solutions to the complex Monge-Ampère and Hessian equations.

The following examples are due to Dinew and Ko\l odziej (see \cite{DK14}).
\skpt
\begin{example}
\begin{enumerate}
\item
Dinew and Ko\l odziej proved in \cite{DK14} that the volume measure $\lambda_{2 n}$ is dominated by capacity. Namely for any $r$ such that $1 < r < \spfrac{m}{n - m}$, there exists a constant $N (r) > 0$ such that for any compact subset $K \subset \Omega$,
\begin{equation} \label{eq:DK}
\lambda_{2 n} (K) \leq N (r) \Capa_m (K,\Omega)^{1+ r}.
\end{equation}
Observe that this estimate is sharp in terms of the exponent when $m < n$. This can be seen by taking $\Omega = \B$ the unit ball and $K:= \bar{\B_s} \subset \B$ the closed ball of radius $s \in{}]0,1[$, since $\Capa_m(\bar{\B}_s,\B) \approx s^{2 (n-m)}$ as $s \to 0$ (see \cite{Lu12}).
When $m=n$ we know that the domination is much more precise (see \cite{ACKPZ09}).

\item
Let $ 0 \leq f \in L^p (\Omega)$ with $p > n \slash m$. Then $ \sfrac{n (p-1)}{p (n - m)} > 1$. By Hölder inequality and inequality \eqref{eq:DK} we obtain: for any $\tau\in{}]1,\sfrac{n (p-1)}{p (n - m)}[$ there exists a constant $M (\tau) > 0$ such that for any compact set $K \subset \Omega$,
\begin{equation} \label{eq:DK2}
\int_K f d \lambda_{2 n} \leq M (\tau) \Vert f\Vert_p \Capa_m (K,\Omega)^{\tau}.
\end{equation}
\end{enumerate}
\end{example}

Theorem~\ref{th:A} will provide us with many new examples. Condition \eqref{eq:capdomination} plays an important role in the following stability result which will be a crucial point in the proof of our theorems (see \cite{EGZ09, GKZ08, Ch16}).

\begin{proposition} \label{prop:stability} Let $\mu$ be a positive Borel measure on $\Omega$ dominated by the $m$\nobreakdash-Hessian capacity with parameters $(A, \tau)$ such that $\tau > 1$.

Then for any $u, v \in \mathcal {SH}_m (\Omega) \cap L^{\infty} (\Omega)$ such that $(dd^c u)^m \wedge \beta^{n - m} \leq \mu$ on $\Omega$ and $\liminf_{\partial \Omega} (u - v) \geq 0$, we have
\begin{equation} \label{eq:stability}
\sup_{\Omega} (v - u)_+ \leq 2 \Vert (v-u)_+\Vert_{1,\mu}^{1 \slash (m+1)} + C \|(v-u)_+\|^\gamma_{1,\mu},
\end{equation}
where $\Vert (v-u)_+ \Vert_{1,\mu}:= \int_{\Omega} (v-u)_+ d \mu$ and
\begin{equation} \label{eq:stbilityconstant}
C:= 1 + \frac{2^{\tau} A^{\sfrac{1}{m}}}{1-2^{1 -\tau}}, \quad \gamma = \gamma (\tau,m):= \frac{\tau - 1}{\tau (m + 1) - m }\cdot
\end{equation}
\end{proposition}
Observe that the most relevant case in applications is when $\Vert (v-u)_+ \Vert_{1,\mu}$ is small. So the right exponent is $\gamma < 1 \slash (m +1)$.
\begin{proof}
The proof uses an idea which goes back to Ko\l odziej (\cite{Kol96}) with some simplifications due to Guedj, Eyssidieux and the second author (see \cite{EGZ09, GKZ08}). It relies on the following estimates: for any $t > 0$, $s > 0$
\begin{equation}\label{eq:Cap-MA}
t^m \Capa_m(\{u< v -s-t\},\Omega)\leq \int_{\{u<v -s\}}(dd^cu)^m\wedge\beta^{n-m}.
\end{equation}
Indeed let $t > 0$, $s > 0$ be fixed and $ w \in \mathcal{SH}_m(\Omega)$ be given such that $-1 \leq w \leq0$. Then
\[\{u-v<-s-t\}\subset\{u-v< t w -s\} \subset\{u-v<-s\} \Subset\Omega.\]
It follows that
\begin{align*}
t^m \int_{\{u-v<-s-t\}}(dd^c w)^m\wedge\beta^{n-m}&\leq\int_{\{u< v - s-t\}}(dd^c(v+tw))^m\wedge\beta^{n-m}\\
&\leq \int_{\{u< v+ tw - s\}}(dd^c(v+tw))^m\wedge\beta^{n-m}.
\end{align*}

On the other hand the comparison principle yields
\begin{align*}
\int_{\{u< v+ tw - s\}}(dd^c(v+tw))^m\wedge\beta^{n-m} &\leq\int_{\{u < v +t w - s\}}(dd^c u)^m\wedge\beta^{n-m}\\
&\leq\int_{\{u< v -s\}}(dd^cu)^m\wedge\beta^{n-m}.
\end{align*}
The last two inequalities imply \eqref{eq:Cap-MA}.

Applying inequality \eqref{eq:Cap-MA} with the parameter $(s \slash 2, s\slash 2)$ instead of $(t,s)$ and taking into account that $u$ is a supersolution, we obtain
\[
\begin{split}
\Capa_m(\{u< v - s \},\Omega) & \leq 2^m s^{-m} \int_{\{u<v - s\slash 2\}}(dd^cu)^m\wedge\beta^{n-m} \\
& \leq 2^{m + 1} s^{-m - 1} \int_\Omega (v-u)_+ d \mu.
\end{split}
\]
Set $s_0:= 2 \Vert (v-u)_+\Vert_{1,\mu}^{1 \slash (m+1)}$. Then for any $s \geq s_0$,
\begin{equation} \label{eq:smallcap}
\Capa_m(\{u< v - s \},\Omega)\leq 1.
\end{equation}

Fix $\varepsilon > 0$ and $s \geq 0$. Then applying inequality \eqref{eq:Cap-MA} with $s_0 + s + \varepsilon$ instead of $s$ and taking into account the fact that $(dd^c u)^m \wedge \beta^{n - m} \leq \mu$ weakly on $\Omega$, we get
\begin{equation} \label{eq:cap-mu}
t^m \Capa_m(\{u< v - s_0 - \varepsilon - s -t\},\Omega)\leq \int_{\{u<v -s_0 - \varepsilon - s\}} d \mu.
\end{equation}
Set $ f (s) = f_\varepsilon (s):=\Capa_m(\{u-v<-s - s_0 -\varepsilon\},\Omega)^{\sfrac{1}{m}}$. By \eqref{eq:smallcap}, we have $f (s) \leq 1$. Hence since $\mu$ is dominated by capacity, it follows that for any $t > 0$ and $s > 0$,
\[
t f (s + t) \leq A^{\sfrac{1}{m}} f (t)^{1 + a}, \quad \text{where } a:= \tau - 1 > 0.
\]
It follows from \cite[Lem.\,2.4]{EGZ09}) that $f (s) = 0$ for any $s\geq S_\infty$ where
\[
S_\infty:=\frac{2 A^{\sfrac{1}{m}} }{1-2^{-a}} [f (0)]^{a},
\]
Thus $v-u\leq s_0 + \varepsilon + S_\infty$ almost everywhere on $\Omega$ and then the inequality holds everywhere on $\Omega$, \ie
\[
\max (v-u)_+ \leq s_ 0 + \varepsilon + \frac{2 A^{\sfrac{1}{m}}}{1-2^{-a}}\, \Capa_m(\{v - u >\varepsilon\},\Omega)^{a}
\]
Applying \eqref{eq:Cap-MA} with $t=\varepsilon$ and $s = 0$ we obtain
\[
\Capa_m(\{v - u > \varepsilon\},\Omega)\leq 2 \varepsilon^{-m-1}\|(v-u)_+\|_{1,\mu}.
\]
As a consequence of the previous estimate, we obtain
\[
\sup_\Omega(v-u)\leq 2 \Vert (v-u)_+\Vert_{1,\mu}^{1 \slash (m+1)} + \varepsilon+ C' \varepsilon^{-a (m+1)}\|(v-u)_+\|^{a}_{1,\mu},
\]
where $C':= \spfrac{2^{a + 1} A^{\sfrac{1}{m}}}{1-2^{-a}}$.
Set $\varepsilon:=\|(v-u)_+\|^\gamma_{1,\mu}$, with
\[
\gamma:=\frac{a}{1 +a(m+1)} = \frac{\tau - 1}{(\tau - 1) (m + 1) + 1}.
\]
Then
\[
\sup_{\Omega}(v-u)_+\leq 2 \Vert (v-u)_+\Vert_{1,\mu}^{1 \slash (m+1)} + C \|(v-u)_+\|^\gamma_{1,\mu},
\]
where $C:= C' + 1 = 1 + \spfrac{2^{a + 1} A^{\sfrac{1}{m}}}{1-2^{-a}} = 1 + \spfrac{2^{\tau} A^{\sfrac{1}{m}}}{1-2^{1 -\tau}}$.
\end{proof}

\section{Subharmonic envelopes and obstacle problems}\label{sec:3}
Here we prove some results that will be used in the proof of Theorem~\ref{th:A}. Since they are of independent interest, we will state them in the most general form and give complete proofs.

\subsection{Subharmonic envelopes }

Let $\Omega \Subset \C^n$ and $h: \Omega \to \R$ is a non positive bounded Borel function and define the corresponding projection:
\begin{equation} \label{eq:subextension}
\tilde h = P_{m,\Omega} (h):= \left(\sup \{v \in \mathcal {SH}_m (\Omega) \sep v \leq h \, \text{in} \, \, \Omega\}\right)^*.
\end{equation}
Observe that we do not need to take the upper semi-continuous regularization if $h$ is upper semi-continuous on $\Omega$. On the other hand, we can easily see that
\[
P_{m,\Omega} (h):= \sup \{v \in \mathcal {SH}_m (\Omega) \sep \, v \leq h \, \text{quasi everywhere on} \, \, \Omega\},
\]
where $v \leq h$ quasi everywhere on $\Omega$ means that the exceptional set where $v \geq h$ has zero $\Capa_m$-capacity.

This is a classical construction in potential theory and has been considered in Complex Analysis first by H.\,Bremermann in \cite{Brem59}, J.\,B.\,Walsh in \cite{Wal69} and also by J.\,Siciak in \cite{Sic81}. Later it has been studied by Bedford and Taylor when solving the Dirichlet problem for the the complex Monge-Ampère equation (\cite{BT76}, \cite{BT82}. In the setting of compact Kähler manifolds it has bee considered R.\,Berman and J.-P.\,Demailly in \cite{BD12} and later in \cite{Ber19}. It has been also considered recently in \cite{GLZ19} in connexion with the supersolution problem for complex Monge-Ampère equations, where a precise estimate of its complex Monge-Ampère measure was given.

We will extend these last results to Hessian equations.

\begin{lemma} \label{lem:projection} Let $\Omega \Subset \C^n$ be a bounded strongly $m$-pseudoconvex domain and $h$ a bounded lower semi-continuous function on $\Omega$. Then the function $\tilde h:= P_{m,\Omega} (h)$ satisfies the following properties:
\begin{enumeratei}
\item\label{lem:projectioni}
$\tilde h \in \mathcal {SH}_m (\Omega) \cap L^{\infty} (\Omega)$, and $\tilde h \leq h$ a.e. on $\Omega$;

\item\label{lem:projectionii}
if $h$ is continuous on $\bar \Omega$, then $\tilde h$ is continuous on $\bar{\Omega}$ and satisfies the following properties
\begin{equation} \label{eq:boundaryvaluesh}
\lim_{\Omega \ni z \to \zeta} \tilde h (z) = h (\zeta), \quad \zeta \in \partial \Omega,
\end{equation}

\item\label{lem:projectioniii}
$ \int_\Omega (\tilde h - h) (dd^c \tilde h)^m \wedge \beta^{n - m} = 0$.
\end{enumeratei}
\end{lemma}

\begin{proof}
Observe that $\min_{\bar \Omega} h \leq \tilde h \leq \max_{\bar \Omega} h$ on $\Omega$.
\begin{enumerate}
\item
Property \eqref{lem:projectioni} follows from the general theory (see \cite{Lu12}).

\item
Property \eqref{lem:projectionii} can be proved using the perturbation method due to J.\,B.\,Walsh (see \cite{Wal69}). Let us recall the argument for completeness.

We first prove that $\tilde h$ satisfies \eqref{eq:boundaryvaluesh} meaning that it has boundary values equal to~$h$ and then it extends as a function on $\bar \Omega$ which is continuous on $\partial \Omega$. Indeed fix $\varepsilon > 0$ and let $h'$ be a $C^2$ approximating function on $\bar \Omega$ such that $h - \varepsilon \leq h' \leq h $ on $\bar \Omega$. Let $\rho$ be the strongly $m$-subharmonic defining function for $\Omega$. Then there exists a constant $A > 0$ such that $u:= A \rho + h'$ is $m$-subharmonic on $\Omega$ and $u \leq h' \leq h$ on $\bar \Omega$. Hence by definition of the envelope, we have $u \leq \tilde h \leq h$ on $\bar \Omega$. Therefore, for any $\zeta \in \partial \Omega$,
\begin{align*}
h (\zeta) - \varepsilon \leq h' (\zeta) & = \lim_{\Omega \ni z \to \zeta } u (z) \\
& \leq \liminf_{\Omega \ni z \to \zeta } \tilde h (z) \leq \limsup_{\Omega \ni z \to \zeta } \tilde h (z) \leq h (\zeta).
\end{align*}
Since $\varepsilon > 0$ is arbitrary, we obtain the identity \eqref{eq:boundaryvaluesh}.
We can then extend $\tilde h$ to $\bar \Omega$ by setting $\tilde h (\zeta) = h (\zeta)$ for $\zeta \in \partial \Omega$. To prove continuity of $\tilde h$ on $\bar \Omega$, we use the perturbation argument of J.\,B.\,Walsh. Fix $\delta > 0$ small enough, $a \in \C^n$ such that $\vert a \vert \leq \delta$ and set $\Omega_a:= (- a) + \Omega$. We define the modulus of continuity of $\tilde{h}$ near the boundary as follows:
\[
\tilde{\kappa}_{\tilde h} (\delta):= \sup \{ \vert \tilde h(z) - \tilde{h} (\zeta)\vert \sep z \in \Omega, \zeta \in \partial \Omega, \vert z - \zeta\vert \leq \delta.
\]
Then, since $\tilde h = h$ is uniformly continuous on $\partial \Omega$, we see that $\lim_{\delta \to 0^+} \tilde{\kappa}_{\tilde h} (\delta) = 0$. By~definition of $\tilde{\kappa}_{\tilde h}$, for any $z \in \Omega \cap \partial \Omega_a$, we have
\[\tilde{ h} (z + a) \leq \tilde h (z) + \tilde{\kappa}_{\tilde{h}} (\delta) \leq \tilde h (z) + \tilde{\kappa}_{\tilde{h}} (\delta) + \kappa_h (\delta),
\]
where $ \kappa_h (\delta)$ is the modulus of continuity of $h$ on $\bar{\Omega}$. Therefore, by the gluing principle, the function defined by
\[
v (z):=\begin{cases}
\max \{ \tilde h (z), \tilde h (z + a) - \tilde{\kappa}_{\tilde{h}} (\delta) - \kappa_h (\delta)\} &\hbox{if } z \in\Omega \cap \Omega_a \\
\tilde h (z) &\hbox{if } z \in\Omega\setminus\Omega_a
\end{cases}
\]
is $m$-subharmonic on $\Omega$ and satisfies $v \leq h$ on $\bar \Omega$. Hence $v \leq \tilde h$ on $\bar \Omega$ and then
\[
\tilde h (z + a) - \tilde{\kappa}_{\tilde{h}} (\delta) - \kappa_h (\delta) \leq \tilde h (z),
\]
for any $z \in \Omega \cap \Omega_a$ with $\vert a\vert \leq \delta$. This proves that $\tilde{h}$ is uniformly continuous on $\bar{\Omega}$.

\item
Property \eqref{lem:projectioniii} follows by a standard balayage argument in potential theory which goes back to Bedford and Taylor for the complex Monge-Ampère equation (\cite{BT76}, \cite{BT82}, see also \cite{GLZ19}).\qedhere
\end{enumerate}
\end{proof}

\begin{remark}
The proof above does not give any information on the modulus of continuity of $\tilde h$ in terms of the modulus of continuity of $h$. In other words we do not know if $\tilde{\kappa}_{\tilde h}$ is comparable to $\kappa_h$.

However if $h$ is $C^2$-smooth on $\bar \Omega$, the function $u:= A \rho + h$, considered in the proof above with $h'=h$, is $m$-subharmonic on $\Omega$, Lipschitz on $\bar{\Omega}$ and satisfies $u \leq \tilde{h} \leq h$ on $\bar{\Omega}$.
Then this implies that $\tilde{\kappa}_{\tilde{h}} (\delta) \leq \kappa_h (\delta) + \kappa_u (\delta) \leq C \kappa_h (\delta)$, where $C > 0$ is a uniform constant.
Therefore, the modulus of continuity of $\tilde h$ satisfies the inequality $\kappa_{\tilde h}(\delta) \leq C' \kappa_h(\delta)$,
where $C' > 0$ is an absolute constant.

This information is not needed here, but it is worth mentioning that this an interesting open problem related to the regularity of solutions to obstacle problems.
We~will come back to this in a subsequent work.
\end{remark}

\Subsection{An obstacle problem}

\begin{theorem} \label{thm:obstacle}
Let $h \in \Cc^2 ({\bar{\Omega}})$. Then $\tilde h:= P_{m,\Omega} h \in \mathcal{SH}_m (\Omega) \cap C^0 (\bar \Omega)$ and its $m$\nobreakdash-Hessian measure satisfies the following inequality:
\begin{equation} \label{eq:BerIneq1}
(dd^c \tilde h)^m \wedge \beta^{n - m} \leq \mathbf{1}_{\{ \tilde h = h\}} \sigma_m^+ (h),
\end{equation}
in the sense of currents on $\Omega$.
\end{theorem}
Here for a function $h \in \Cc^2 ({\bar{\Omega}})$, we set
\[
\sigma_m^+ (h):= \mathbf{1}_G \, \sigma_m (h),
\]
pointwise on $\Omega$, where $G$ is the set of points $z \in \Omega$ such that $dd^c h (z) \in \Theta_m$ \ie the $(1,1)$-form $dd^c h (z)$ is $m$-positive (see Definition \ref{def:mpositive}).

\begin{proof} To prove \eqref{eq:BerIneq1}, we proceed as in \cite{GLZ19}, using an idea which goes back to R.\,Berman \cite{Ber19}. Thanks to the property \eqref{lem:projectionii} of Lemma \ref{lem:projection}, it is enough to prove that
\begin{equation} \label{eq:BerIneq2}
(dd^c \tilde h)^m \wedge \beta^{n - m} \leq \sigma_m^+ (h),
\end{equation}
in the sense of currents on $\Omega$. We proceed in two steps.

\begin{enumerate}
\item
Assume first that $\Omega$ is smooth strongly $m$-pseudoconvex and $h \in \Cc^2 (\bar \Omega)$ and consider the following Dirichlet problem for the complex $m$-Hessian equation depending on the parameter $j \in \N$,
\begin{equation} \label{eq:BerEqu}
(dd^c u)^m \wedge \beta^{n - m} = e^{j (u-h)} \sigma_m^+ (h), \quad u = h\text{ in }\partial \Omega.
\end{equation}
By \cite{Lu13}, for each $j\!\in\!\N$, there exists a unique continuous solution $u_j\!\in\!\mathcal {SH}_m (\Omega) \cap\Cc^0 (\Omega)$ to this problem (see also \cite{Ch16}).

Our goal is to prove that the sequence $(u_j)_{j\in \N}$ increases to $ \tilde h$ uniformly on $\bar{\Omega}$. We argue as in \cite{GLZ19} with obvious modifications. Recall $h$ is $C^2$ in $\bar{\Omega}$. Then by~definition~$h$ is a viscosity supersolution to the Dirichlet problem \eqref{eq:BerEqu}. Moreover, by~Proposition~\ref{prop: viscosity vs potential general case}, $u_j$ is a viscosity subsolution to the Dirichlet problem \eqref{eq:BerEqu}. By the viscosity comparison principle Theorem \ref{thm: viscosity comparison principle}, we conclude that $u_j \leq h$ in $ \Omega$ since $u_j = h$ on~$\partial \Omega$.

Therefore, the pluripotential comparison principle (Proposition \ref{prop:comparison principle}) implies that $(u_j)$ is an increasing sequence.
On the other hand, by Theorem \ref{thm:boundedsubsolution} there exists a bounded $m$-subharmonic function $\psi$ on $\Omega$ which is a solution to the complex Hessian equation
\[
\sigma_m (\psi) = e^{\psi - h} \sigma_m^+ (h)
\]
with $\psi = h$ on $\partial \Omega$.
Moreover, for any $j \in \N$, one can easily check that the function defined by the formula
\[
\psi_j:= (1- 1 \slash j) \tilde h + (1 \slash j) ( \psi - m \log j)
\]
is a (pluripotential) subsolution to the equation \eqref{eq:BerEqu}, since $ \tilde h \leq h$ on $\Omega$. Hence by Proposition \ref{prop:comparison principle} we have $\psi_j \leq u_j$ on $\Omega$.

Summarizing, we have proved that for any $j \in \N$, $\psi_j \leq u_j \leq \tilde h $ on $\Omega$.
Therefore,
\[
0 \leq \tilde h - u_j \leq \tilde h - \psi_j = (1 \slash j) ( \tilde h - \psi + m \log j)\quad\text{on $\Omega$ for any $j \in \N^*$}.
\]
This proves that $u_j$ converges to $\tilde h$ uniformly on $\Omega$. Then since $u_j \leq h$ on $\Omega$, taking the limit as $j \to + \infty$ in \eqref{eq:BerEqu} we obtain inequality \eqref{eq:BerIneq2} by the continuity of the Hessian operators for uniform convergence (see \cite{Lu12}).

\item
For the general case of a bounded $m$-hyperconvex domain, we approximate $\Omega$ by an increasing sequence $(\Omega_j)_{j \in \N}$ of smooth strongly $m$-pseudoconvex domains such that for any $j \in \N$, $\Omega_{j + 1} \subset \Omega_j$ and $\Omega = \cup_{j \in \N} \Omega_j$.
Then it is easy to see that the sequence $(P_{m,\Omega} h_j)$ decreases to $P_{m,\Omega} h$ on $\Omega$ (see \cite{GLZ19}). Thus the result follows from the previous case by the continuity of the Hessian operator for monotone sequences.\qedhere
\end{enumerate}
\end{proof}

It is worth mentioning that these envelopes have been considered by several authors in the context of compact Kähler manifolds. When $h$ is $C^2$ it was proved recently that $P (h)$ is $C^{1,1}$ (see \cite{ChZh17}, \cite{T18}, \cite{Ber19}) and equality holds in \eqref{eq:BerIneq1}, which means that $P(h)$ is a solution to an obstacle problem (see \cite{BD12}).

We can address a similar question.

\begin{Question}
Is it true that $\tilde h$ is $C^{1,1}$ locally on $\Omega$ when $h$ is $C^2$ on $\bar \Omega$? Is there equality in \eqref{eq:BerIneq1}?
\end{Question}

\begin{corollary} \label{cor:smoothing} Let $\Omega \Subset \C^n$ be a strongly $m$-pseudoconvex domain. Let $u \in \mathcal {SH}_m (\Omega)$ a negative $m$-subharmonic function. Then there exists a decreasing sequence $(u_j)$ of continuous $m$-subharmonic functions on $\Omega$ with boundary values $0$ which converges pointwise to $u$ on $\Omega$.
\end{corollary}

\begin{proof} We can assume that $u$ is bounded on $\Omega$ and extend it as a semi-continuous function on $\bar \Omega$. Let $(h_j)_{j \in \N}$ be a decreasing sequence of smooth functions in a neighbourhood of $\bar \Omega$ which converges to $u$ in $\bar \Omega$.
For each $j \in \N$, consider the $m$-subharmonic envelope
$v_j:= P_{\Omega} h_j$ on $\Omega$ and set $u_j:= \max \{v_j, j \rho\}$ on $\Omega$, where $\rho$ is a continuous $m$-subharmonic defining function for $\Omega$. Then by Lemma \ref{lem:projection}, $(u_j)$ is a decreasing sequence of continuous $m$-subharmonic functions on $\Omega$ which converges to $u$ on $\Omega$.
\end{proof}

Applying the smoothing method of Richberg it is possible to construct a decreasing sequence of smooth $m$-subharmonic functions on $\Omega$ which converges to $u$ in $ \Omega$ (see \cite{P14}).

\section{Hessian measures of Hölder continuous potentials}\label{sec:4}

In this section we will prove two important results which will be used in the proof of the main theorems stated in the introduction.

\subsection{Hessian mass estimates near the boundary}

Here we prove a comparison inequality which seems to be new even in the case of a complex Monge-Ampère measure.

\begin{lemma} \label{lem:ComparisonIneq}
Let $\Omega \Subset \C^n$ be a bounded strongly $m$-pseudoconvex domain and $\varphi \in \mathcal{SH}_m(\Omega)\cap \Cc^{\alpha} (\bar{\Omega})$ ($0 < \alpha \leq 1$) with $\varphi \equiv 0$ on $\partial \Omega$. Then for any Borel set $K \subset \Omega$, we have
\[
\int_K (dd^c\varphi)^m\wedge\beta^{n-m} \leq L^m \left[\delta_K (\partial \Omega)\right]^{ m \alpha} \, \Capa_m (K,\Omega),
\]
where
\[
\delta_K (\partial \Omega):= \sup_{z \in K} \dist (z, \partial \Omega)
\]
and $L > 0$ is the Hölder norm of $\varphi$.
\end{lemma}
The constant $\delta_K (\partial \Omega)$ is the Hausdorff distance of $K$ to the boundary in the sense that $\delta_K (\partial \Omega) \leq \varepsilon $ means that $K$ is contained in the $\varepsilon$-neighbourhood of $\partial \Omega$.

The relevant point here is that the estimate takes care of the behaviour at the boundary. It shows in particular that if the volume of the compact set is fixed, the capacity tends to $+ \infty$ when the compact set approaches the boundary at a rate controlled by the Hausdorff distance of the compact to the boundary.
\begin{proof}
By inner regularity, we can assume that $K \subset \Omega$ is compact.
Since $\varphi$ is Hölder continuous on $\bar \Omega$, we have
$\varphi (\zeta) - \varphi (z) \leq L \vert \zeta - z\vert^{\alpha}$ for any $\zeta \in \partial \Omega$ and any $z \in \Omega$.

Fix a compact set $K \subset \Omega$. Since $\varphi = 0$ in $\partial \Omega$, it follows that for any $z \in K$,
\[
- \varphi (z) \leq \kappa \left[\dist (z,\partial \Omega)\right]^{\alpha} \leq L \left[\delta_K (\partial \Omega)\right]^{\alpha} =:a.
\]
Therefore, the function $v:= a^{-1} \varphi \in \mathcal{SH}_m (\Omega)$ and $ v \leq 0$ on $\Omega$ and $v \geq - 1$ in $K$. Fix $\varepsilon >0$ and let $u_K$ be the relative extremal $m$-subharmonic function of $(K,\Omega)$. Then $ K \subset \{ (1 + \varepsilon) u_K^* < v\} \cup \{u_K < u_K^*\}$. Since the set $\{u_K < u_K^*\}$ has zero $m$-capacity (see \cite{Lu12}), it follows from the comparison principle that for any $\varepsilon > 0$,
\begin{align*}
\int_K (dd^c v)^m \wedge \beta^{n - m} &\leq \int_{\{ (1 + \varepsilon) u_K^* < v\}} (dd^c v)^m \wedge \beta^{n - m} \\
&\leq (1+ \varepsilon)^m \int_{\{ (1 + \varepsilon) u_K^* < v\}} (dd^c u_K^*)^m \wedge \beta^{n - m} \\
& \leq (1+ \varepsilon)^m \Capa_m (K,\Omega).
\end{align*}
The last inequality follows from \eqref{eq:cap}.
The estimate of the lemma follows by letting $\varepsilon \to 0$.
\end{proof}

\subsection{Hölder continuity of Hessian measures}

In order to prove the Hölder continuous subsolution theorem we need an additional argument following an idea which goes back to \cite{DDGKPZ15} and used in a systematic way in \cite{N18a} (see also \cite{KN19}).

Given a continuous function $g \in \Cc^0 (\partial \Omega)$ and a real number $R > 0$, we denote by $\mathcal{E}^g_m(\Omega,R)$ the convex set of bounded $m$-subharmonic functions $v $ on $\Omega$ such that $v = g$ on $\partial \Omega$ normalized by the mass condition $\int_{\Omega} (dd^c v)^m \wedge \beta^{n -m} \leq R$.

In order to prove Theorem~\ref{th:B}, we will need the following Lemma:
\begin{lemma}\label{lem:ModC}
Let $\varphi\in \mathcal{E}^0_m(\Omega)\cap \Cc^{\alpha} (\overline\Omega)$ ($ 0 < \alpha \leq 1$) and $g \in \Cc^0 (\partial \Omega)$ and $R > 0$.
Then for any $1 \leq k\leq m$, there exists $C_k = C (k,m,\Omega,R) >0$ such that for every $u,v\in \mathcal{E}^g_m(\Omega,R)$
\begin{equation} \label{eq:MocEst}
\int_{\Omega}|u-v|(dd^c\varphi)^k\wedge\beta^{n-k}\leq C_k \left[\Vert u-v \Vert_1\right]^{\alpha_k},
\end{equation}
provided that $\Vert u-v \Vert_1:= \int_{\Omega} \vert u - v\vert d \lambda_{2 n} \leq 1$, where $\alpha_k:= \alpha^k 2^{-k}$.
\end{lemma}
\begin{proof}
Recall the following notation for the complex Hessian measure of $\varphi$:
\[
\sigma_k{(\varphi)}:= (dd^c\varphi)^k\wedge\beta^{n-k} \ \ \ 1 \leq k\leq m.
\]

Observe that for any $\varepsilon > 0,$ $ u_\varepsilon:= \max \{u - \varepsilon, v \} \in \mathcal{E}^g_m(\Omega)$ and $u_\varepsilon = v$ near the boundary $\partial \Omega$. By the the comparison principle, this implies that
$u_\varepsilon \in \mathcal{E}^g_m(\Omega,R)$. Therefore, replacing $u$ by $u_\varepsilon$, we can assume that $u \geq v$ on $\Omega$ and $u = v$ near the boundary $\partial \Omega$. Then the inequality \eqref{eq:MocEst} will follow from this case since $\vert u - v \vert = ( \max\{u,v\} - u) + (\max\{u,v\} - v)$.

On the other hand by approximation on the support $S$ of $u-v$ which is compact, we can assume that $u$ and $ v $ are smooth on a neighbourhood of $S$.

We will argue by induction on $k$. For $k=0$, the inequality is obviously satisfied with $C_0 = 1$. Assume that the inequality holds for some integer $k$ such that $0 \leq k<m$ \ie
\begin{equation}\label{eq:Hyprec}
\int_{\Omega}( u-v) \sigma_k{(\varphi)}\leq C_{k} \left[\Vert u-v \Vert_1\right]^{\alpha_k}.
\end{equation}
We will show that there exists $C_{k+1} > 0$ such that
\[
\int_{\Omega}(u-v) \sigma_{k+1}(\varphi)\leq C_{k+1} \left[\Vert u-v \Vert_1\right]^{\alpha_{k+1}}.
\]
We will approximate $\varphi$ by smooth functions. We first extend $\varphi$ as a Hölder continuous function on $\C^n$. Indeed recall that for any $z, \zeta \in \bar{\Omega}$, we have $\varphi (z) \leq \varphi (\zeta) + \kappa \vert z - \zeta\vert^\alpha $. Then it is easy to see that the following function
\begin{equation} \label{eq:Holderextension}
\bar{\varphi} (z):= \sup \{\varphi (\zeta) - \kappa \vert z - \zeta\vert^\alpha \sep \zeta \in \bar{\Omega}\}, \, \, z \in \C^n.
\end{equation}
is Hölder continuous of order $\alpha$ on $\C^n$ and $ \bar{\varphi} = \varphi$ on $\Omega$. For simplicity, we will denote this extension by $\varphi$. Then we denote by $\varphi_{\delta}$ the smooth approximants of $\varphi $ on $\C^n$, obtained by Formula \eqref{eq:regularization}. By Lemma \ref{lem:Poisson-Jensen} for $\delta\in{}]0,\delta_0[$, $\varphi_\delta \in \mathcal{SH}_m (\Omega_{\delta})\cap\Cc^{\infty}(\C^n)$.

To prove the required estimate, we write
\[
\int_{\Omega}(u-v) (dd^c\varphi)^{k+1} \wedge\beta^{n-k-1} = A + B,
\]
where
\[
A:=\int_{\Omega}(u-v)dd^c \varphi_{\delta}\wedge(dd^c\varphi)^k\wedge\beta^{n-k-1},
\]
and
\[
B := \int_{\Omega}(u-v) dd^c(\varphi_{\delta}-\varphi)\wedge(dd^c\varphi)^k\wedge\beta^{n-k-1}.
\]
We estimate each term separately.
Fix $\delta\in{}]0,\delta_0[$. Since $\varphi$ is Hölder continuous, we have $\vert \varphi_{\delta} (z) - \varphi (z)\vert \leq \kappa \delta^{\alpha}$ for any $z \in \Omega$. Moreover, differentiating the formula \eqref{eq:reg}, we obtain for $z \in \Omega$,
\begin{equation} \label{eq:estimat2}
dd^c \varphi_{\delta} (z) \leq M_1\, \frac{ \kappa \delta^{\alpha}}{\delta^2}\, \beta,
\end{equation}
where $M_1 > 0$ depends only on $\Omega$ and $\chi$. By \eqref{eq:estimat2} and \eqref{eq:Hyprec}, we have
\begin{equation} \label{eq:estimationA}
|A|\leq M_1\, \frac{\kappa \delta^{\alpha}}{\delta^2} \int_{\Omega}(u-v) \sigma_{k}(\varphi)
\leq M_1 C_k \kappa \delta^{\alpha - 2} C_{k} \left[\Vert u-v \Vert_1\right]^{\alpha_k}.
\end{equation}
To estimate $B$, observe that, since $ u - v = 0$ near the boundary, we can integrate by parts to get the following formula
\[
B = \int_{\Omega} (\varphi-\varphi_{\delta}) dd^c (u-v) \wedge(dd^c\varphi)^k\wedge\beta^{n-k-1},
\]
and then
\[
\vert B\vert \leq \int_{\Omega} \vert \varphi_{\delta}-\varphi\vert dd^c (u+v) \wedge(dd^c\varphi)^k\wedge\beta^{n-k-1}.
\]
Therefore, since $\vert \varphi_{\delta}-\varphi \vert \leq \kappa \delta^{\alpha}$ on $\Omega$, it follows that
\begin{equation} \label{eq:estimationB}
|B| \leq( I_k(u,\varphi) + I_k(v,\varphi) ) \, \kappa \delta^{\alpha},
\end{equation}
where $I_k (u,\varphi):= \int_{\Omega} dd^c u \wedge(dd^c\varphi)^k\wedge\beta^{n-k-1}$. Observe that by Lemma \ref{lem:Cegrell} and the normalization mass condition, there exists a constant $d(m,n) > 0$ such that for any~$k$ such that $1 \leq k \leq m$, $I_k(u,\varphi) + I_k(v,\varphi) \leq d (m,n)$. Combining this with the inequalities \eqref{eq:estimationA} and \eqref{eq:estimationB}, we obtain for $\delta\in{}]0,\delta_0[$,
\[
\int_{\Omega}(u-v) \sigma_{k +1}(\varphi)\leq M_1 C_k \kappa \delta^{\alpha- 2} [\|u-v\|_1]^{\alpha_k}+ d(m,n) \kappa \delta^{\alpha}.
\]
Since $ [\|u-v\|_1 \leq 1$., we can take $\delta= \delta_0 \sqrt{ [\|u-v\|_1]^{\alpha_k}} < \delta_0$ in the last inequality to obtain
\begin{align*}
\int_{\Omega}(u-v) \sigma_{k+1}(\varphi) &\leq (M_1 C_k + d (m,n))\, \kappa \left(\sqrt{ [\|u-v\|_1]^{\alpha_k}}\right)^{\alpha} \\
& = C_{k + 1} [\|u-v\|_1]^{\alpha_{k+ 1}},
\end{align*}
where $\alpha_{k + 1}:= \alpha_k (\alpha \slash 2)$.
\end{proof}

It is an open problem to know the precise modulus of continuity of the Hessian measure $\sigma_m (\varphi)$ acting on the space of normalized $m$-subharmonic potentials $\mathcal{E}^g_m(\Omega,R) \subset L^1 (\Omega,\lambda_{2 n})$. We do not know either if the lemma is true when the total mass of the Hessian measure $\sigma_m (\varphi) $ on $\Omega$ is infinite.

\section{Proofs of the main results}\label{sec:5}

In this section we will give the proofs of Theorem~\ref{th:A} and Theorem~\ref{th:B} stated in the introduction using the previous results.

\subsection{ Proof of Theorem~\ref{th:A}}

For the proof of Theorem~\ref{th:A}, we will use the same idea as \cite{KN19}. However, since our measure does not have compact support, we need to use the control on the behaviour of the mass of the $m$-Hessian of the subsolution close to the boundary, given by Lemma \ref{lem:ComparisonIneq}.

\begin{proof}
We extend $\varphi$ as a Hölder continuous function on the whole of $\C^n$ with the same exponent and denote by $\varphi$ the extension (see \eqref{eq:Holderextension}).
Then denote by $\varphi_{\delta}$ ($0 < \delta < \delta_0$) the smooth approximants of $\varphi $ on $\Omega$ defined as usual by Formula \eqref{eq:regularization}. Then we know that $\varphi_\delta \in \mathcal{SH}_m (\Omega_{\delta})\cap\Cc^{\infty}(\C^n)$.

We consider the $m$-subharmonic envelope of $\varphi_\delta$ on $\Omega$ defined by the formula
\[
\psi_\delta:= \sup \{\psi \in \mathcal{SH}_m (\Omega) \sep \psi \leq \varphi_\delta \, \, \, \hbox{on} \, \, \, \Omega \}\cdot
\]
It follows from Lemma \ref{lem:projection} that $\psi_\delta \in \mathcal{SH}_m (\Omega)$ and $\psi_\delta \leq \varphi_\delta $ on $\Omega$. Fix $\delta \in{}]0, \delta _0[$ and a compact set $K \subset \Omega_\delta$, and consider the set
\[
E:=\{3 \kappa \delta^\alpha u_K^*+ \psi_\delta<\varphi- 2 \kappa \delta^{\alpha}\} \subset \Omega.
\]
Since $\varphi $ is Hölder continuous on $\bar \Omega$, we have $\varphi - \kappa \delta^\alpha\leq \varphi_{\delta} \leq \varphi + \kappa \delta^{\alpha}$ on $\Omega$ and then $\varphi - \kappa \delta^\alpha \leq \psi_\delta \leq \varphi_{\delta} \leq \varphi (z) + \kappa \delta^\alpha$ on $\Omega$.
Therefore, $\liminf_{z \to \partial \Omega} (\psi_\delta - \varphi + \kappa \delta^\alpha ) \geq 0$, and then $E \Subset \Omega$. By the comparison principle, we conclude that
\begin{equation}\label{eq:fundmentalestimate}
\begin{split}
\int_{E}(dd^c\varphi)^m\wedge&\beta^{n-m} \leq \int_{E}(dd^c(3 \kappa \delta^\alpha u_K^* + \psi_{\delta}))^m\wedge\beta^{n-m} \\
& \leq 3 \kappa L \delta^\alpha \int_{E}(dd^c(u_K^* + \psi_{\delta}))^m\wedge\beta^{n-m}+\int_{E}(dd^c\psi_{\delta})^m\wedge\beta^{n-m},
\end{split}
\end{equation}
where $L:= \max_{0 \leq j \leq m - 1} (3 \kappa \delta_0^\alpha)^j$.

Observe that $-1 + \varphi - \kappa \delta^\alpha \leq u_K^* + \psi_\delta \leq \varphi + \kappa \delta^\alpha$ on $\Omega$, hence $\vert u_K^* + \psi_\delta\vert \leq \sup_{\Omega} \vert \varphi\vert + 1 + \kappa \, \delta_0^\alpha =: M_0$ on $\Omega$. Therefore, from inequality \eqref{eq:fundmentalestimate}, it follows that
\begin{equation} \label{eq:finalestimate1}
\int_{E}(dd^c\varphi)^m\wedge\beta^{n-m} \leq 3 \kappa \delta^\alpha L M_0^m \Capa_m (E,\Omega) + \int_{E}(dd^c\psi_{\delta})^m\wedge\beta^{n-m}.
\end{equation}
Since $\varphi$ is Hölder continuous on $\bar \Omega$, we have
\begin{equation}\label{eq:1a}
dd^c\varphi_{\delta}\leq\frac{M_1 \kappa \delta^\alpha}{\delta^{2}}\, \beta \quad\text{ on } \Omega,
\end{equation}
where $M_1 > 0$ is a uniform constant depending only on $\Omega$. Hence by Theorem \ref{thm:obstacle}, we have
\begin{equation}\label{eq:1b}
(dd^c\psi_{\delta})^m\wedge\beta^{n-m} \leq (\sigma_m (\varphi_{\delta}))_+ \leq\frac{M_1^m \kappa^m \delta^{m \alpha}}{\delta^{2 m}}\, \beta^n,
\end{equation}
in the sense of currents on $ \Omega$. Therefore,
\[
\int_{E}(dd^c\psi_{\delta})^m\wedge\beta^{n-m}\leq M_1^{m} \kappa^m \delta^{ m (\alpha -2)} \lambda_{2 n} (E).
\]
From this estimate and the inequalities \eqref{eq:finalestimate1} and \eqref{eq:1b}, we deduce that
\begin{equation} \label{eq:finalestimate2}
\int_{E}(dd^c\varphi)^m\wedge\beta^{n-m} \leq 3 \kappa \delta^{\alpha} L M_0^m \Capa_m (E,\Omega)+ M_1^{m} \kappa^m \delta^{(\alpha -2)m} \lambda_{2 n} (E).
\end{equation}
By the volume-capacity comparison inequality \eqref{eq:DK}, it follows that for any fixed $r$ with $1 < r < \spfrac{m}{n - m}$, there exists a constant $N (r) > 0$ such that
\begin{equation} \label{eq:volumeestimate}
\lambda_{2 n}(E) \leq N (r) [\Capa_m(E,\Omega)]^{1 + r}.
\end{equation}
Since $ E \subset\{u_K^* < - 1 \slash 3\}$, by the comparison principle we deduce the following inequality:
\begin{equation}\label{eq:2}
\Capa_m (E,\Omega) \leq 3^{m} \Capa_m (K,\Omega).
\end{equation}
Since $K \setminus \{u_K < u_K^*\} \subset E$ and $K \cap \{u_K < u_K^*\}$ has zero capacity, it follows that $\int_K (dd^c\varphi)^m\wedge\beta^{n-m} \leq \int_E (dd^c\varphi)^m\wedge\beta^{n-m}$.

Therefore, if we set $c_m (\cdot):= \Capa_m (\cdot,\Omega)$, we finally deduce from \eqref{eq:finalestimate2}, \eqref{eq:volumeestimate} and \eqref{eq:2} that for a fixed $\delta\in{}]0,\delta_0[$ and any compact set $K \subset \Omega_\delta$, we have
\begin{equation}
\int_K (dd^c\varphi)^m\wedge\beta^{n-m} \leq C_0\kappa \delta^{\alpha} c_m (K) + C_1 \kappa^m \delta^{(\alpha -2) m } [c_m(K)]^{1+ r}.
\end{equation}
where $C_0:= 3^{m + 1} L M_0^m$ and $C_1:= M_1^{m} 3^{m r} N (r)$. By inner regularity of the capacity, we deduce that the previous estimate holds for any Borel subset $B \subset \Omega_\delta$, \ie \begin{equation} \label{eq:estimate1}
\int_B (dd^c\varphi)^m\wedge\beta^{n-m} \leq C_0 \kappa \delta^{\alpha} c_m (B) + C_1 \kappa^{m \alpha } \delta^{(\alpha -2) m} [c_m(B)]^{1+r}.
\end{equation}

Let $K \subset \Omega$ be any fixed compact set and $\delta\in{}]0,\delta_0[$. Then
\[
\int_K (dd^c\varphi)^m\wedge\beta^{n-m} = \int_{K \cap \Omega_\delta} (dd^c\varphi)^m\wedge\beta^{n-m} + \int_{K \setminus \Omega_\delta} (dd^c\varphi)^m\wedge\beta^{n-m}.
\]
We will estimate each term separately. By \eqref{eq:estimate1} the first term is estimated easily:
\[
\int_{K \cap \Omega_\delta} (dd^c\varphi)^m\wedge\beta^{n-m} \leq C_0 \kappa \delta^{\alpha} c_m (K) + C_1 \kappa^{m \alpha}\delta^{-2m + m \alpha} [c_m(K)]^{1+ r}.
\]
To estimate the second term we apply Lemma \ref{lem:ComparisonIneq} for the Borel set $ B:= K \setminus \Omega_\delta$. Since $\delta_B (\partial \Omega) \leq \delta$ we get
\[
\int_{K \setminus \Omega_\delta} (dd^c\varphi)^m\wedge\beta^{n-m} \leq \kappa^m \delta^{m \alpha} c_m (K).
\]
Therefore, we obtain the following estimate. For any $\delta\in{}]0,\delta_0[$ and any compact set $K \subset \Omega$, we have
\begin{multline} \label{eq:fundIneq}
\int_{K}(dd^c\varphi)^m\wedge\beta^{n-m}\\[-5pt]
\leq C_0 \kappa \delta^{\alpha} c_m (K) + C_1 \kappa^m \delta^{(\alpha -2) m } [c_m(K)]^{1+ r} + \kappa^m \delta^{m \alpha} c_m (K).
\end{multline}
We want to optimize the right hand side of \eqref{eq:fundIneq} by taking $\delta:= [c_m (K)]^{\spfrac{r}{(2-\alpha) m + \alpha}}$. Observe that if $ \delta_K(\partial \Omega) \leq [c_m (K)]^{\spfrac{r}{(2-\alpha) m + \alpha}}$, then by Lemma \ref{lem:ComparisonIneq} we get
\begin{equation} \label{eq:finalIneq1}
\int_{K}(dd^c\varphi)^m\wedge\beta^{n-m} \leq \kappa^m [c_m (K)]^{1 + \spfrac{m \alpha r}{(2-\alpha) m + \alpha}}.
\end{equation}
Now assume that $ [c_m (K)]^{\spfrac{r}{(2-\alpha) m + \alpha}} < \delta_K (\partial \Omega) \leq \delta_0$. Then we can take $\delta:= [c_m (K)]^{\spfrac{r}{(2-\alpha) m + \alpha}}$ in inequality \eqref{eq:fundIneq} and get
\begin{equation}\label{eq:finalIneq2}
\int_{K}(dd^c\varphi)^m\wedge\beta^{n-m} \leq (C_0 \kappa + C_1 \kappa^m + \kappa^m) \, [c_m (K)]^{1 + \spfrac{\alpha r}{(2-\alpha) m + \alpha}}.
\end{equation}
Combining inequalities \eqref{eq:finalIneq1} ) and \eqref{eq:finalIneq2}, we obtain the estimate of the theorem with the constant $A$ given by the following formula:
\begin{equation} \label{eq:finalConst}
A:= C_0 \kappa + C_1 \kappa^m + \kappa^m.\qedhere
\end{equation}
\end{proof}

\subsection{Proof of Theorem~\ref{th:B}}

Now we are ready to prove Theorem~\ref{th:B} from the introduction using Theorem~\ref{th:A} and Lemma \ref{lem:ModC}.

\begin{proof} According to Theorem \ref{thm:boundedsubsolution}, we know that there is a unique function $u\in \mathcal{SH}_m(\Omega)\cap L^{\infty}(\overline{\Omega}) $ such that
\[
(dd^c u)^m\wedge\beta^{n-m}=\mu,
\]
in the weak sense on $\Omega$ and $ u=g $ on $\partial\Omega$. To complete the proof we need to show that~$u$ is Hölder continuous up to the boundary.

For $\delta\in{}]0,\delta_0[$ and denote as before by $u_{\delta}(z)$ the $\delta$-regularization of~$u$. Recall that ${u}_{\delta}$ is $m$-subharmonic on $\Omega_\delta$. We construct a global $m$-subharmonic function $\tilde{u}_{\delta}$ which is close to ${u}_{\delta}$ on $\Omega_\delta$.

By \cite{Ch16} there exists a continuous maximal $m$-subharmonic function $w \in \mathcal{SH} (\Omega) \cap C^{\alpha} (\bar \Omega)$ such that $w = g$ on $\partial \Omega$. Then $v: = w + \varphi \in \mathcal{SH} (\Omega) \cap \Cc^{\alpha} (\bar \Omega)$ is a subsolution to the Dirichlet problem \eqref{eq:DirPb} such that $v = g$ on $\partial \Omega$. Hence $v \leq u \leq w$. To prove Hölder continuity on $\Omega$, it's enough by Lemma \ref{lem:sup-mean} to estimate $u_\delta:= u \star \chi_\delta$
on $\Omega_\delta$.

We claim that there exists a constant $\kappa > 0$ such that for $z\in\partial\Omega_{\delta}$, we have $ u (z) \geq {u}_{\delta}(z) - \kappa \delta^{\alpha} $. Indeed fix $z\in\partial\Omega_{\delta}$. Then there exists $\zeta \in \partial \Omega$ such that $\vert z -\zeta \vert = \delta$. Since $v \leq u \leq w$ on $\Omega$ and they are equal on $\partial \Omega$, it follows that
\begin{align*}
u_{\delta} (z) & \leq w_\delta (z) \leq w (z) + \kappa_w \delta^\alpha \\
& \leq w (\zeta) + 2 \kappa_w \delta^\alpha = v (\zeta) + 2 \kappa_w \delta^\alpha \\
& \leq v (z) +( \kappa_v + 2 \kappa_w) \delta^\alpha \\
& \leq u (z) + \kappa \delta^\alpha,
\end{align*}
where $\kappa:= \kappa_v + 2 \kappa_w$ and $ \kappa_v$ (\resp $\kappa_w$) is the Hölder constant of $v$ (\resp $w$). This proves our claim.

Therefore, the following function
\[
\tilde{u}_{\delta}:= \begin{cases}
\max\{{u}_{\delta} - \kappa \delta^{\alpha},u \} &\text{on } \Omega_{\delta}, \\
u &\text{on } \Omega\setminus\Omega_{\delta}
\end{cases}
\]
is $m$-subharmonic and bounded on $\Omega$ and satisfies
\[
0 \leq \tilde{u}_{\delta}(z) - u (z) =(u_\delta (z) - u (z) - \kappa \delta^{\alpha})_+ \leq u_\delta (z) - u (z)\quad \text{for }z\in \Omega_{\delta}
\]
and $\tilde{u}_{\delta}(z) - u (z) = 0$ on $\Omega \setminus \Omega_\delta$. Moreover, since $ \tilde{u}_{\delta} \geq u$ on $ \Omega$ and $\tilde{u}_{\delta} = u$ on $\Omega \setminus \Omega_\delta$, Corollary \ref{coro:comparison principle} implies that
\[
\int_{\Omega_\delta}(dd^c\tilde{u}_{\delta})^m\wedge\beta^{n-m}\leq\int_{\Omega_\delta}(dd^c u)^m\wedge\beta^{n-m}.
\]
Hence for any $\delta\in{}]0,\delta_0[$, we have
\[
\int_{\Omega}(dd^c\tilde{u}_{\delta})^m\wedge\beta^{n-m} \leq \mu (\Omega) < \infty.
\]
Since $\tilde u_\delta \geq u$ on $\Omega$, Proposition \ref{prop:stability} implies that for any $\gamma$ such that
\[
0<\gamma < \gamma (m,n,\alpha):= \frac{m \alpha}{ m (m + 1) \alpha + 2 (n-m)},
\]
there exists a constant $D_\gamma > 0$ such that any $\delta\in{}]0,\delta_0[$,
\begin{equation} \label{eq2}
\sup_{\Omega}(\tilde{u}_{\delta}-u) \leq D_{\gamma} \left(\int_{\Omega}(\tilde{u}_{\delta}-u) d\mu\right)^{\gamma}.
\end{equation}

On the other hand, since $\mu\leq (dd^c\varphi)^m\wedge\beta^{n-m}$ on $\Omega$, it follows from Theorem~\ref{th:A} that we can apply Lemma \ref{lem:ModC} and get for $\delta\in{}]0,\delta_0[$,
\begin{align*}
\int_{\Omega}(\tilde{u}_{\delta}-u)d\mu & \leq C_m \left(\int_{\Omega}(\tilde{u}_{\delta}-u)(z)d\lambda(z)\right)^{\alpha_m} \nonumber \\
&\leq C_m\left(\int_{\Omega_{\delta}}({u}_{\delta}(z)-u(z)d\lambda(z)\right)^{\alpha_m}.
\end{align*}
By Lemma \ref{lem:Poisson-Jensen}, the previous inequality implies that
\begin{equation} \label{eq3}
\int_{\Omega}(\tilde{u}_{\delta}-u)d\mu \leq C_m(B \|\Delta u\|_{\Omega}\delta^2)^{\alpha_m}.
\end{equation}
Since $ \max \{{u}_{\delta}- \kappa \delta^{\alpha},u\}-u\leq {u}_{\delta}-u$ on $\Omega$, it follows from the equations \eqref{eq2} and \eqref{eq3} that
\begin{align*}
\sup_{\Omega_{\delta}}( {u}_{\delta}-u)&\leq\sup_{\Omega}(\tilde{u}_{\delta}-u)+ \kappa \delta^{\alpha} \\
&\leq C_m^{\gamma} D_\gamma \left(B\|\Delta u\|_{\Omega}\delta^2\right)^{\gamma\alpha_m} + \kappa \delta^{\alpha}.
\end{align*}
By Lemma \ref{lem:Cegrell}, we have $\|\Delta u\|_{\Omega} \leq c_{m,n} \mu (\Omega)^{1 \slash m} < + \infty$. Then for $\delta\in{}]0,\delta_0[$,
\[
\sup_{\Omega_{\delta}}({u}_{\delta}-u)\leq C' (m,n,\alpha) \delta^{2 \gamma \alpha_m}.
\]
Since $2 \gamma \alpha_m < \alpha$, it follows from Lemma \ref{lem:sup-mean} that for $\delta\in{}]0,\delta_0[$ and $z \in \Omega_\delta$,
\[
\sup_{B (z,\delta)} u \leq u (z) + C'' (m,n,\alpha) \delta^{2 \gamma \alpha_m},
\]
where $C''(m,n,\alpha) > 0$ is a positive constant which can be made explicitly using the proof in \cite{GKZ08}.
This proves Hölder continuity of $u$ on $\bar \Omega$.
\end{proof}
